%% TRANSCRIBED (PDF 174-191 = printed IV-29..IV-46)
%% NOTE: printed pages IV-32 and IV-33 are BLANK placeholder pages in the
%% working PDF (Taylor & Francis ebook defect; local PDF 177-178 contain only
%% the publisher logo).  They were recovered from the textually identical
%% 1989 Addison-Wesley photo-reprint of the 1968 Benjamin edition (public
%% McGill scan, math.mcgill.ca/darmon/courses/17-18/sato-tate/serre-mcgill.pdf;
%% same 203-page pagination), and were checked to continue verbatim into
%% IV-31 and IV-34 of the working copy -- the same defect as the one noted
%% in ch3_s2.tex (printed III-6/III-7).
%% TODO-VERIFY(p177,p178): the IV-32/IV-33 passages below follow the McGill
%% scan render (checked at 200-500 dpi), not the working PDF (blank pages).

\apphead{Local results}{LOCAL RESULTS}

In what follows, $K$ denotes a field which is complete with respect to a
discrete valuation $v$; we denote by $O_K$ (resp. by $k$) the ring of
integers (resp. the residue field) of $K$; we assume that $k$ is perfect
and of characteristic $p\neq 0$.

Let $E$ be an elliptic curve over $K$ and let $\ell$ be a prime number
different from the characteristic of $K$. Let $T_\ell$ and $V_\ell$ be
the corresponding Galois modules; we denote by $G_\ell$ the image of
$\Gal(K_s/K)$ in $\Aut(T_\ell)$, and by $I_\ell$ the inertia subgroup of
$G_\ell$. The Lie algebras $\mathfrak{g}_\ell=\mathrm{Lie}(G_\ell)$,
$\mathfrak{i}_\ell=\mathrm{Lie}(I_\ell)$ are subalgebras of $\End(V_\ell)$
and we will determine them under suitable assumptions on $K$ and $v$;
note that, since $I_\ell$ is an invariant subgroup of $G_\ell$, its Lie
algebra $\mathfrak{i}_\ell$ is an \emph{ideal} of $\mathfrak{g}_\ell$.

If $j=j(E)$ is the modular invariant of $E$ (cf.~\S\ref{sec:IV.1.1}), we
consider the cases $v(j)<0$ and $v(j)\geq 0$ separately.

\subsechead{A.1}{The case $v(j)<0$}

In this section we assume that the modular invariant $j$ of the elliptic
curve $E$ has a pole, i.e. that $v(j)<0$.

\subsubsechead{A.1.1}{The elliptic curves of Tate}

Let $q$ be an element of $K$ with $v(q)>0$, and let $\Gamma_q$ be the
discrete subgroup of $K^*$ generated by $q$. Then, by Tate's theory of
ultrametric theta functions (unpublished -- but see Morikawa, Nagoya
Math.~Journ., 1962), there is an elliptic curve $E_q$ defined over $K$
with the property that, for any finite extension $K'$ of $K$, the
analytic group $K'^*/\Gamma_q$ is isomorphic to the group $E_q(K')$ of
points of $E_q$ with values in $K'$. The equation defining $E_q$ can be
written in the form
\[
y^2 + xy = x^3 - b_2x - b_3\;,
\]
with
\[
b_2 = 5\sum_{n\ge 1} n^3q^n/(1-q^n)
\quad\text{and}\quad
b_3 = \sum_{n\ge 1} (7n^5+5n^3)q^n/12(1-q^n)\;,
\]
these series converging in $K$. The modular invariant $j(q)$ of $E_q$ is
given by the usual formula
\[
j(q) = \frac{(1+48b_2)^3}{q\prod_{n\ge 1}(1-q^n)^{24}}
     = \frac{1}{q} + 744 + 196884q + \dotsb\;,
\]
a series with integral coefficients. The function field of $E_q$ consists
of the fractions $F/G$, where $F$ and $G$ are Laurent series
\[
F = \sum_{-\infty}^{+\infty} a_nz^n\,,\qquad
G = \sum_{-\infty}^{+\infty} b_nz^n
\]
with coefficients in $K$, converging for all values of $z\neq 0,\infty$,
and such that $F(qz)/G(qz) = F(z)/G(z)$.

Since the modular invariant $j$ of the given elliptic curve $E$ is such
that $v(j)<0$, and since the series for $j(q)$ has integral coefficients,
one can choose $q$ so that $j=j(q)$. The elliptic curves $E$ and $E_q$
become then isomorphic over a finite extension of $K$ (which can be taken
to be of degree $2$). Hence, after possibly replacing $K$ by a finite
extension, \emph{we may assume that} $E=E_q$.

\subsubsechead{A.1.2}{An exact sequence}

We conserve the notation of \S\ref{sec:IV.A.1.1}. Let $E_n$ be the kernel
of multiplication by $\ell^n$ in $K_s^*/\Gamma_q$. If $\mu_n$ is the
group of $\ell^n$-th roots of unity in $K_s$, we have an injection
$\mu_n\longrightarrow E_n$. On the other hand, if $z\in E_n$, we have
$z^{\ell^n}\in\Gamma_q$, and hence there exists an integer $c$ such that
$z^{\ell^n}=q^c$. If we associate to $z$ the image of $c$ in
$\Z/\ell^n\Z$, we obtain a homomorphism of $E_n$ into $\Z/\ell^n\Z$, and
the resulting sequence
\begin{equation}
\tag{1}\label{eq:IV.A.1.2.1}
0 \longrightarrow \mu_n \longrightarrow E_n
   \longrightarrow \Z/\ell^n\Z \longrightarrow 0
\end{equation}
is an exact sequence of $\Gal(K_s/K)$-modules, $\Gal(K_s/K)$ acting
trivially on $\Z/\ell^n\Z$. Passing to the limit, we obtain an exact
sequence of Galois modules
\begin{equation}
\tag{2}\label{eq:IV.A.1.2.2}
0 \longrightarrow T_\ell(\mu) \longrightarrow T_\ell(E_q)
   \longrightarrow \Z_\ell \longrightarrow 0
\end{equation}
where $\Gal(K_s/K)$ acts trivially on $\Z_\ell$. Tensoring with
$\Q_\ell$, we obtain the exact sequence
\begin{equation}
\tag{3}\label{eq:IV.A.1.2.3}
0 \longrightarrow V_\ell(\mu) \longrightarrow V_\ell(E_q)
   \longrightarrow \Q_\ell \longrightarrow 0\;.
\end{equation}

We now show that this sequence of $\Gal(K_s/K)$-modules does not split.
To do this we introduce an invariant $x$ which belongs to the group
$\varprojlim H^1(G,\mu_n)$, where $G=\Gal(K_s/K)$. Let $d$ be the
coboundary homomorphism:
\[
H^0(G,\Z/\ell^n\Z) \longrightarrow H^1(G,\mu_n)
\]
%% p177 (IV-32) follows the McGill scan; see header note.
with respect to the exact sequence \eqref{eq:IV.A.1.2.1} and let
$x_n=d(1)$. The invariant $x$ is the element of
$\varprojlim H^1(G,\mu_n)$ defined by the family $(x_n)$, $n\geq 1$.

\begin{proposition}
\begin{alist}
\item The isomorphism
$\delta\colon K^*/K^{*\ell^n}\longrightarrow H^1(G,\mu_n)$ of Kummer
theory transforms the class of $q$ mod $K^{*\ell^n}$ into $x_n$.
\item The element $x$ is of infinite order.
\end{alist}
\end{proposition}

(Recall that $\delta$ is induced by the coboundary map relative to the
exact sequence
\[
1 \longrightarrow \mu_n \longrightarrow \bar{K}^*
  \xrightarrow{\ \ell^n\ } \bar{K}^* \longrightarrow 1.)
\]

Assertion (a) is proved by an easy computation. To prove (b), note that
the valuation $v$ defines a homomorphism
\[
f_n\colon K^*/K^{*\ell^n} \longrightarrow \Z/\ell^n\Z\;,
\]
and hence a homomorphism
\[
f\colon \varprojlim K^*/K^{*\ell^n} \longrightarrow \Z_\ell\;.
\]
If we identify $x$ with the corresponding element of
$\varprojlim K^*/K^{*\ell^n}$, as in (a), we have $f(x)=v(q)$, hence $x$
is of infinite order.

\begin{corollary}
The sequence \eqref{eq:IV.A.1.2.3} does not split.
\end{corollary}

Assume it does, i.e. there is a $G$-subspace $X$ of $V_\ell(E_q)$ which
is mapped isomorphically onto $\Q_\ell$. Let $X_T = T_\ell(E_q)\cap X$.
The image of $X_T$ in $\Z_\ell$ is $\ell^N\Z_\ell$, for some $N\geq 0$.
It is then easy to see that $\ell^Nx=0$, and this contradicts the fact
that $x$ is of infinite order.

\subsubsechead{A.1.3}{Determination of $g_l$ and $i_l$}
%% p178-179 (IV-33) up to ``It follows'' follows the McGill scan; the
%% continuation on IV-34 is from the working PDF.

We keep the notation of \S\ref{sec:IV.A.1.1} and \S\ref{sec:IV.A.1.2}.
If $X$ is a one-dimensional subspace of $V_\ell = V_\ell(E)$, let
$\mathfrak{r}_X$ denote the subalgebra of $\End(V_\ell)$ consisting of
those endomorphisms $u$ for which $u(V_\ell)\subset X$, and let
$\mathfrak{n}_X$ be the subalgebra of $\mathfrak{r}_X$ formed by those
$u\in\mathfrak{r}_X$ with $u(X)=0$.

\begin{theorem}
\begin{alist}
\item If $k$ is algebraically closed and $\ell\neq p$, then there is a
one-dimensional subspace $X$ of $V_\ell$ such that
$\mathfrak{g}_\ell=\mathfrak{n}_X$.
\item If $k$ is algebraically closed and $\ell=p$, then there is a
one-dimensional subspace $X$ of $V_\ell$ such that
$\mathfrak{g}_\ell=\mathfrak{r}_X$.
\item If $k$ is finite, then $\mathfrak{g}_\ell=\mathfrak{r}_X$ for some
one-dimensional subspace $X$ of $V_\ell$, and
$\mathfrak{i}_\ell=\mathfrak{n}_X$ (resp.
$\mathfrak{i}_\ell=\mathfrak{r}_X$) if $\ell\neq p$ (resp. $\ell=p$).
\end{alist}
\end{theorem}

\begin{proof}
Note first that, since $\mathfrak{g}_\ell$ and $\mathfrak{i}_\ell$ are
invariant under finite extension of $K$, we may assume that $E=E_q$.

(a) In this case, $K$ contains the $\ell^n$-th roots of unity, hence
$\Gal(K_s/K)$ acts trivially on $T_\ell(\mu)$. Consequently, there is a
basis $e_1, e_2$ of $T_\ell(E)$ such that, for all
$\sigma\in\Gal(K_s/K)$, we have $\sigma(e_1)=e_1$,
$\sigma(e_2)=a(\sigma)e_1+e_2$ with $a(\sigma)\in\Z_\ell$. Moreover, the
homomorphism $\sigma\longmapsto a(\sigma)$ cannot be trivial since the
sequence \eqref{eq:IV.A.1.2.3} does not split. It follows that
$\Ima(a)$ is an open subgroup of $\Z_\ell$, and hence that
$\mathfrak{g}_\ell=\mathfrak{n}_X$ with $X=V_\ell(\mu)$.

(b) Since $\ell=p$, we must have $\mathrm{char}(K)=0$ as
$\ell\neq\mathrm{char}(K)$. In this case, the action of $\Gal(\bar K/K)$
on $V_\ell(\mu)$ is by means of the character $\chi_\ell$
(cf.~chap.~\ref{chap:I}, \S\ref{sec:I.1.2}) which is of infinite order.
It follows that $\mathfrak{g}_\ell=\mathfrak{r}_X$, where
$X=V_\ell(\mu)$; in fact, $\mathfrak{g}_\ell\supset\mathfrak{n}_X$ since
the sequence \eqref{eq:IV.A.1.2.3} does not split, and we cannot have
$\mathfrak{g}_\ell=\mathfrak{n}_X$.

(c) Since $k$ is finite, the action of $\Gal(K_s/K)$ on $T_\ell(\mu)$ is
not trivial nor even of finite order. Hence, the argument used in (b)
shows that $\mathfrak{g}_\ell=\mathfrak{r}_X$, where $X=V_\ell(\mu)$.
Applying (a) to the completion of the maximal unramified extension of
$K$, we see that $\mathfrak{i}_\ell=\mathfrak{n}_X$ if $\ell\neq p$, and
that $\mathfrak{i}_\ell=\mathfrak{r}_X$ if $\ell=p$.
\end{proof}

\runhead{Exercise}

In case (a), shows that $\Ima(a)=\ell^n\Z_\ell$, where $\ell^n$ is the
highest power of $\ell$ which divides $v(q)=-v(j)$.

\subsubsechead{A.1.4}{Application to isogenies}

Here, we assume that $k$ is \emph{finite} and $K$ \emph{is of
characteristic} $0$ (i.e. $K$ is a finite extension of $\Q_p$).

\begin{theorem}
Let $q, q'\in K^*$ with $v(q)$ and $v(q')>0$. Let $E=E_q$ and $E'=E_{q'}$
be the corresponding elliptic curves over $K$. Then the following are
equivalent:
\begin{alist}
\item[(1)] $E_q$ is $K$-isogenous to $E_{q'}$.
\item[(2)] There are integers $A, B\geq 1$ such that $q^A=q'^B$.
\item[(3)] $V_p(E)$ and $V_p(E')$ are isomorphic as
$\Gal(\bar K/K)$-modules.
\end{alist}
\end{theorem}

\begin{proof}
(2)~$\Rightarrow$~(1). It suffices to show that $E_q$ and $E_{q^A}$ are
isogenous over $K$. But every meromorphic function $F/G$ invariant under
multiplication by $q$ is invariant under multiplication by $q^A$; hence
the function field of $E_q$ is contained in the function field of
$E_{q^A}$, i.e., $E_q$ and $E_{q^A}$ are isogenous.

(1)~$\Rightarrow$~(3). Trivial.

(3)~$\Rightarrow$~(2). Choose an isomorphism $\varphi$ of $V_p(E)$ onto
$V_p(E')$. Since $V_p(\mu)$ is the only one-dimensional subspace of
$V_p(E)$ (resp. $V_p(E')$) stable by $G=\Gal(\bar K/K)$, $\varphi$ maps
$V_p(\mu)$ into itself. Moreover, after multiplying $\varphi$ by an
homothety, we may suppose that $\varphi$ maps $T_p(E)$ into $T_p(E')$.
We then have a commutative diagram:
\begin{equation}
\tag{4}\label{eq:IV.A.1.4.4}
\begin{array}{ccccccccc}
0 &\longrightarrow& T_p(\mu) &\longrightarrow& T_p(E)
  &\longrightarrow& \Z_p &\longrightarrow& 0\\
& & \rho\,\big\downarrow & & \varphi\,\big\downarrow
  & & \big\downarrow\,\sigma & & \\
0 &\longrightarrow& T_p(\mu) &\longrightarrow& T_p(E')
  &\longrightarrow& \Z_p &\longrightarrow& 0
\end{array}
\end{equation}
where $\rho$ (resp. $\sigma$) is the multiplication by a $p$-adic
integer $r$ (resp. $s$). If $x, x'$ are the elements of
$\varprojlim H^1(G,\mu_n)$ associated to $E$ and $E'$
(cf.~\S\ref{sec:IV.A.1.2}), the commutativity of
\eqref{eq:IV.A.1.4.4} shows that
\[
rx = sx'\;.
\]
But the valuation $v$ yields a homomorphism of
$\varprojlim H^1(G,\mu_n)=\varprojlim K^*/K^{*p^n}$ into $\Z_p$, and we
have seen that the image of $x$ is $v(q)$, and the image of $x'$ is
$v(q')$. Hence
\[
rv(q) = sv(q')\;.
\]
We will now show that the element
\[
z = q^{v(q')}/q'^{v(q)}
\]
is a root of unity. First of all, the image of $z$ in
$\varprojlim K^*/K^{*p^n}$ is a $p^a$-th root of unity; in fact, this
image is
\[
v(q')x - v(q)x'\;,
\]
and multiplying by $s$, we find $0$ in virtue of the above formulae
(note that $\varprojlim K^*/K^{*p^n}$ is a $\Z_p$-module, hence
multiplication by $s$ makes sense). We then use the fact that the kernel
of
\[
K^* \longrightarrow \varprojlim K^*/K^{*p^n}
\]
is $k^*$ (viewed, as usual as a subgroup of $K^*$). To see this, one
decomposes $K^*$ as a product $\Z\times k^*\times U^1$, where $U^1$ is
the group of units congruent to $1$. The functor
\[
A \longmapsto \varprojlim A/A^{p^n}
\]
transforms $\Z$ into $\Z_p$, kills $k^*$ and leaves $U^1$ unchanged,
since $U^1$ is a finitely generated $\Z_p$-module. Hence, we have
$z\in k^*$, and $z$ is a root of unity. This implies (1), q.e.d.
\end{proof}

\runhead{Remark}

The equivalence (1)~$\Longleftrightarrow$~(2) was remarked by Tate. It
is true without any hypothesis on $K$.

\runhead{Exercise}

Show that the hypothesis ``$k$ is finite'' may be replaced by ``$k$ is
algebraic over $\F_p$.''

\subsubsechead{A.1.5}{Existence of transvections in the inertia group}

Let $E$ be the elliptic curve $E_q$ (cf.~\S\ref{sec:IV.A.1.1}), let
$\widetilde{G}_\ell$ be the image of $\Gal(K_s/K)$ in
$\Aut(T_\ell/\ell T_\ell)$, and let $\widetilde{I}_\ell$ be the inertia
subgroup of $\widetilde{G}_\ell$. We assume that $v$ is normalized, i.e.
that $v(K^*)=\Z$.

\begin{proposition}
If $\ell$ does not divide $v(q)$, then $\widetilde{I}_\ell$ contains a
transvection, i.e. an element whose matrix is
$\begin{pmatrix} 1 & 1 \\ 0 & 1 \end{pmatrix}$ for a suitable
$\F_\ell$-basis of $T_\ell/\ell T_\ell$.
\end{proposition}

\begin{proof}
After possibly replacing $K$ by a larger field, we can suppose that the
residue field $k$ is algebraically closed, and that $K$ contains the
$\ell$-th roots of unity. In fact, if $\ell\neq p$, this last condition
is implied by the first; if $\ell=p$, we must adjoin these roots; but
the degree of the extension thus obtained divides $\ell-1$, hence is
prime to $\ell$, and the valuation of $q$ remains prime to $\ell$. This
being said, the hypothesis on $v(q)$ shows that $q^{1/\ell}$ is not in
$K$. Thus there is an automorphism $s\in\Gal(K_s/K)$ such that
$s(q^{1/\ell})=zq^{1/\ell}$, with $z\neq 1$. Then $z$ is a primitive
$\ell$-th root of unity, and $z, q^{1/\ell}$ form a basis of $T_\ell$
modulo $\ell T_\ell$. Since $s(z)=z$, we see that the image of $s$ in
$\widetilde{G}_\ell=\widetilde{I}_\ell$ is the required transvection.
\end{proof}

\subsechead{A.2}{The case $v(j)\geq 0$}

In this section we assume that the modular invariant $j$ of the elliptic
curve $E$ is integral, i.e. that $v(j)\geq 0$. Hence, after possibly
replacing $K$ by a finite extension, we may assume that $E$
\emph{has good reduction} (cf.~\S\ref{sec:IV.1.2}). We also assume that
$K$ \emph{is of characteristic zero}.

\subsubsechead{A.2.1}{The case $l\neq p$}

Suppose that $\ell\neq p$. Since $E$ has good reduction, the module
$T_\ell$ can be identified with the Tate module $T_\ell(\widetilde{E})$
of the reduced curve $\widetilde{E}$, cf.~\S\ref{sec:IV.1.3}. Hence the
inertia algebra $\mathfrak{i}_\ell$ is $0$. If the residue field $k$ is
finite, the group $G_\ell$ is topologically generated by the Frobenius
$F_\ell$. Hence, in this case,
$\mathfrak{g}_\ell=\mathrm{Lie}(G_\ell)$ is a one-dimensional subalgebra
of $\End(V_\ell)$.

\subsubsechead{A.2.2}{The case $l=p$ with good reduction of height 2}

Here we assume that the reduced curve $\widetilde{E}$ is of height $2$;
recall that, if $A$ is an abelian variety defined over a field of
characteristic $p$, its height can be defined as the integer $h$ for
which $p^h$ is the inseparable part of the degree of the homothety
``multiplication by $p$.'' An elliptic curve is of height $2$ if and
only if its \emph{Hasse invariant} (cf.~Deuring~\cite{deuring}) is $0$.
Since $E$ has good reduction, it defines an \emph{abelian scheme} $E_v$
over $O_K$, hence a $p$-\emph{divisible group} $E(p)$ over $O_K$
(cf.~Tate~\cite{tate-pdiv}, 2.1 -- see also \cite{serre-pdiv}, \S1,
Ex.~2). The Tate module of $E(p)$ can be identified with $T_p$. The
connected component $E(p)^0$ of $E(p)$ coincides with the \emph{formal
group} (over $O_K$) attached to $E_v$; the height of $\widetilde{E}$ is
precisely the height of this formal group (in the usual sense). In our
case, we have $E(p)=E(p)^0$ since the height is assumed to be $2$.

\begin{theorem}
One has $\mathfrak{g}_p=\mathfrak{i}_p$. This Lie algebra is either
$\End(V_p)$ or a non-split Cartan subalgebra of $\End(V_p)$.
\end{theorem}

(Recall that a non-split Cartan subalgebra of $\End(V_p)$ is a
commutative subalgebra of rank $2$ with respect to which $V_p$ is
irreducible. It is given by a quadratic subfield of $\End(V_p)$.)

\begin{proof}
The Lie algebra $\mathfrak{g}_p$ has the property that
$\mathfrak{g}_pz=V_p$ for any non zero element $z$ of $V_p$
(cf.~\cite{serre-galpdiv}, p.~128, Prop.~8). In particular, $V_p$ is an
irreducible $\mathfrak{g}_p$-module; its commuting algebra is either a
field of degree $2$ (which is then necessarily equal to
$\mathfrak{g}_p$) or the field $\Q_p$, in which case $\mathfrak{g}_p$ is
\emph{a priori} $sl_2$ or $gl_2$. But $\mathfrak{g}_p\neq sl_2$ since
$\Lambda^2V_p$ is canonically isomorphic to $V_p(\mu)$, and the action
of $\Gal(\bar K/K)$ on $V_p(\mu)$ is by means of the character $\chi_p$,
which is of infinite order (indeed, no finite extension of $K$ can
contain all $p^n$-th roots of unity, $n=1,2,\dotsc$). Hence the Lie
algebra $\mathfrak{g}_p$ is either $\End(V_p)$ or a non split Cartan
subalgebra of $\End(V_p)$. Since the above applies to the completion of
the maximal unramified extension of $K$, we have the same alternative
for $\mathfrak{i}_p$. Moreover, $\mathfrak{i}_p$ is contained in
$\mathfrak{g}_p$. We have \emph{a priori} three possibilities:
\begin{alist}
\item $\mathfrak{i}_p=\mathfrak{g}_p=\End(V_p)$.
\item $\mathfrak{i}_p=\mathfrak{g}_p$ is a non split Cartan subalgebra
of $\End(V_p)$.
\item $\mathfrak{i}_p$ is a Cartan subalgebra and
$\mathfrak{g}_p=\End(V_p)$.
\end{alist}
However, $\mathfrak{i}_p$ is an \emph{ideal} of $\mathfrak{g}_p$. Hence,
(c) is impossible, and this proves the theorem.
\end{proof}

\begin{remarks}
\begin{nlist}
\item By a theorem of Tate (\cite{tate-pdiv}, \S4, cor.~1 to th.~4), the
algebra $\mathfrak{g}_p$ is a Cartan subalgebra of $\End(V_p)$ if and
only if $E(p)$ has ``formal complex multiplication,'' i.e. if and only
if the ring of endomorphisms of $E(p)$, over a suitable extension of
$K$, is of rank $2$ over $\Z_p$. There exist elliptic curves without
complex multiplication (in the algebraic sense) whose $p$-completion
$E(p)$ have formal complex multiplication.
\item Suppose that $\mathfrak{g}_p$ is a Cartan subalgebra of
$\End(V_p)$, and let $H=\mathfrak{g}_p\cap\Aut(V_p)$ be the
corresponding Cartan subgroup of $\Aut(V_p)$. If $N$ is the normalizer
of $H$ in $\Aut(V_p)$, then one knows that $N/H$ is cyclic of order
$2$. Since $G_p\subset N$, it follows that $G_p$ is commutative if and
only if $G_p\subset H$. The case $G_p\subset H$ corresponds to the case
where the formal complex multiplication of $E(p)$ is defined over $K$,
and the case $G_p\not\subset H$ corresponds to the case where this
formal multiplication is defined over a quadratic extension of $K$.
\item Suppose that $G_p$ is commutative, and that the residue field $k$
is \emph{finite}. Let $F$ be the quadratic field of formal complex
multiplication (i.e. $\mathfrak{g}_p$ itself, viewed as an associative
subalgebra of $\End(V_p)$). If $U_F$ denotes the group of units of $F$,
the action of $\Gal(\bar K/K)$ on $V_p$ is given by a homomorphism
\[
\varphi\colon \Gal(\bar K/K) \longrightarrow U_F
\]
By local class field theory, we may identify the inertia group of
$\Gal(\bar K/K)^{\mathrm{ab}}$ with the group $U_K$ of units of $K$.
Hence the restriction $\varphi_I$ of $\varphi$ to the inertia group is a
\emph{homomorphism of} $U_K$ \emph{into} $U_F$. To determine
$\varphi_I$, we first remark that the action of $\End(E(p))$ on the
tangent space to $E(p)$ defines an \emph{embedding} of $F$ into $K$. For
that embedding, one has (compare with chap.~\ref{chap:III},
\S\ref{sec:III.A.4})
\[
\varphi_I(x) = N_{K/F}(x^{-1}), \quad\text{for all } x\in U_K
\]
Indeed, by a result of Lubin (Ann.~of Math.~85, 1967), there is a formal
group $E'$ which is $K$-isogenous to $E(p)$, and has for ring of
endomorphisms the ring of integers of $F$. But then, if $E''$ is a
Lubin--Tate group over $K$ (cf.~Lubin--Tate~\cite{lubin-tate}), the
formal groups $E'$ and $E''$ are isomorphic over the completion of the
maximal unramified extension of $K$ (cf.~Lubin~\cite{lubin},
th.~4.3.2). Hence to prove the formula ($*$), we may assume that $E(p)$
is a Lubin--Tate group, in which case the formula ($*$) follows from the
main result of \cite{lubin-tate}.
\end{nlist}
\end{remarks}

\subsubsechead{A.2.3}{Auxiliary results on abelian varieties}

Let $A$ and $B$ be two abelian varieties over $K$, with good reduction,
so that the associated $p$-divisible groups $A(p)$ and $B(p)$ are defined
(these are $p$-divisible groups \emph{over the ring} $O_K$,
cf.~Tate~\cite{tate-pdiv}). Let $\widetilde{A}$ and $\widetilde{B}$
(resp. $\widetilde{A}(p)$ and $\widetilde{B}(p)$) be the reductions of
$A$ and $B$ (resp. of $A(p)$ and $B(p)$).

\begin{theorem}[1]\label{thm:IV.A.2.3.1}
Let $\widetilde{f}\colon \widetilde{A}\to\widetilde{B}$ be a morphism of
abelian varieties, and let $\widetilde{f}(p)$ be the corresponding
morphism of $\widetilde{A}(p)$ into $\widetilde{B}(p)$. Assume there is
a morphism $f(p)\colon A(p)\to B(p)$ whose reduction is
$\widetilde{f}(p)$. Then, there is a morphism
$f\colon A\to B$ whose reduction is $\widetilde{f}$.
\end{theorem}

A proof of this ``lifting'' theorem has been given by Tate in a Seminar
(Woods Hole, 1964), but has not yet been published; a different proof
has been given by W.~Messing (L.~N.~264, 1972).

\begin{theorem}[2]\label{thm:IV.A.2.3.2}
Assume $T_p(A)$ is a direct sum of $\Z_p$-modules of rank $1$ invariant
under the action of $\Gal(\bar K/K)$. Then every endomorphism of
$\widetilde{A}$ lifts to an endomorphism of $A$, i.e., the reduction
homomorphism $\End(A)\to\End(\widetilde{A})$ is surjective (and hence
bijective, since it is known to be injective).
\end{theorem}

Using theorem~\ref{thm:IV.A.2.3.1}, one sees that it is enough to show
that any endomorphism of $\widetilde{A}(p)$ can be lifted to an
endomorphism of $A(p)$. But the assumption made on $T_p$ implies
(cf.~Tate~\cite{tate-pdiv}, 4.2) that $\widetilde{A}(p)$ is a product of
$p$-divisible groups of height $1$. Hence we are reduced to proving the
following elementary result:

\begin{lemma}
Let $H_1, H_2$ be two $p$-divisible groups over $O_K$, both of height
one. Then the reduction map
$\Hom(H_1,H_2)\to\Hom(\widetilde{H}_1,\widetilde{H}_2)$ is bijective.
\end{lemma}

\begin{proof}
This is clear if both $H_1$ and $H_2$ are \'etale. If both are not
\'etale, their duals are \'etale and we are reduced to the previous
case. If one of them is \'etale, and the other is not, one checks
readily that
$\Hom(H_1,H_2)=\Hom(\widetilde{H}_1,\widetilde{H}_2)=0$.
\end{proof}

\begin{corollary}
Assume:
\begin{rlist}
\item $V_p(A)$ is a direct sum of one-dimensional subspaces stable under
$\Gal(\bar K/K)$.
\item The residue field $k$ of $K$ is finite.
\end{rlist}
Then $A$ is isogenous to a product of abelian varieties of (CM) type
(in the sense of Shimura--Taniyama~\cite{shimura-taniyama},
cf.~also chap.~\ref{chap:II}, \S\ref{sec:II.2.8}).
\end{corollary}

\begin{proof}
Assumption (i) implies that $T_p(A)$ contains a lattice $T'$ which is a
direct sum of free $\Z_p$-modules of rank $1$ stable under
$\Gal(\bar K/K)$. One can find an isogeny $A_1\to A$ such that
$T_p(A_1)$ is mapped onto $T'$. This means that, after replacing $A$ by
an isogenous variety, we may apply Th.~\ref{thm:IV.A.2.3.2} to $A$,
i.e. $\End(A)\to\End(\widetilde{A})$ is an isomorphism.
But, since $k$ is finite, it follows from a result of
Tate~\cite{tate-endo} that $\Q\tens\End(\widetilde{A})$ contains a
semi-simple commutative $\Q$-subalgebra $\Lambda$ of rank $2\dim(A)$
(this is not explicitly stated in \cite{tate-endo}, but follows easily
from its ``Main Theorem''). Hence, the same is true for
$\Q\tens\End(A)$. If we now write $\Lambda$ as a product of
commutative fields $\Lambda_a$, one sees that $A$ is isogenous to a
product $\prod_a A_a$, where $A_a$ has complex multiplication of type
$\Lambda_a$, q.e.d.
\end{proof}

\subsubsechead{A.2.4}{The case $l=p$ with good reduction of height 1}

In this section, we assume that the reduced curve $\widetilde{E}$ is of
height $1$ i.e. that its Hasse invariant is $\neq 0$
(cf.~Deuring~\cite{deuring}). The connected component
$E_1=E(p)^0$ of the $p$-divisible group $E(p)$ attached to $E$
(cf.~Tate~\cite{tate-pdiv}) is then a formal group of height $1$. Since
$E(p)$ is an extension of $E_1$ by an \'etale group, we obtain an exact
sequence of $\Gal(\bar K/K)$-modules
\begin{equation}
\tag{$*$}\label{eq:IV.A.2.4.star}
0 \longrightarrow X \longrightarrow V_p \longrightarrow Y
   \longrightarrow 0\;,
\end{equation}
where $X$ corresponds to the Tate module of $E_1$, and $Y$ to the points
of order a power of $p$ of $\widetilde{E}$,

\begin{theorem}
Suppose that the residue field $k$ is finite. Then the following
statements are equivalent:
\begin{alist}
\item[(a)] The elliptic curve $E$ has complex multiplication over $K$.
\item[(a')] The elliptic curve $E$ has complex multiplication over an
extension of $K$.
\item[(b)] There exists a one-dimensional subspace $D$ of $V_p$, which
is a supplementary subspace of $X$, and is stable under the action of
$G_p$.
\item[(b')] There exists a one-dimensional subspace $D$ of $V_p$ which
is a supplementary subspace of $X$, and is stable under the action of
$\mathfrak{g}_p=\mathrm{Lie}(G_p)$.
\end{alist}
\end{theorem}

\begin{proof}
If $D$ is stable under the action of $G_p$, it is also stable under the
action of its Lie algebra $\mathfrak{g}_p$, hence
(b)~$\Rightarrow$~(b'). Conversely, if $D$ is stable under
$\mathfrak{g}_p$, its transforms by $G_p$ are in finite number; a
standard mean value argument then shows that the sequence
\eqref{eq:IV.A.2.4.star} splits, hence (b')~$\Rightarrow$~(b). The
implication (b)~$\Rightarrow$~(a) (the only non-trivial one) follows
from the corollary to theorem~\ref{thm:IV.A.2.3.2} of
\S\ref{sec:IV.A.2.3}. Conversely, if $E$ has complex multiplication by
an imaginary quadratic field $F$, the group $\Gal(\bar K/K)$ acts on
$V_p$ through $F\tens\Q_p$ (see chap.~\ref{chap:II},
\S\ref{sec:II.2.8}) and this action is thus semi-simple. Consequently,
the exact sequence \eqref{eq:IV.A.2.4.star} splits; this shows that
(a)~$\Rightarrow$~(b), hence also that (a')~$\Rightarrow$~(b'). Since
(a)~$\Rightarrow$~(a') is trivial, the theorem is proved.
\end{proof}

\begin{corollary}[1]\label{cor:IV.A.2.4.1}
If $E$ has no complex multiplication, $\mathfrak{g}_p$ is the Borel
subalgebra $\mathfrak{b}_X$ of $\End(V_p)$ formed by those
$u\in\End(V_p)$ such that $u(X)\subset X$; the inertia algebra
$\mathfrak{i}_p$ is the subalgebra $\mathfrak{r}_X$ of
$\mathfrak{b}_X$ formed by those $u\in\End(V_p)$ such that
$u(V_p)\subset X$.
\end{corollary}

Let $\chi_X$ and $\chi_Y$ be the characters of $\Gal(\bar K/K)$ defined
by the one-dimensional modules $X$ and $Y$. Since $k$ is finite,
$\chi_Y$ is of infinite order. If $\chi$ is the character defined by the
action of $\Gal(\bar K/K)$ on $V_p(\mu)$, the isomorphisms
\[
X\tens Y \cong \Lambda^2V_p \cong V_p(\mu)
\]
show that $\chi_X\chi_Y=\chi$. Hence the restriction of $\chi_X$ and
$\chi_X\chi_Y^{-1}$ to the inertia subgroup of $\Gal(\bar K/K)$ are of
infinite order. This shows first that $\mathfrak{g}_p$ is either
$\mathfrak{b}_X$ or a Cartan subalgebra of $\mathfrak{b}_X$; since the
second case would imply (b'), it is impossible, hence
$\mathfrak{g}_p=\mathfrak{b}_X$. Similarly, one sees first that
$\mathfrak{i}_p$ is contained in $\mathfrak{r}_X$, then that its action
on $X$ is non trivial; since it is an ideal in
$\mathfrak{g}_p=\mathfrak{b}_X$, these properties imply
$\mathfrak{i}_p=\mathfrak{r}_X$.

\runhead{Remark}

The above result is given in \cite{serre-lieg}, p.~245, Th.~1, but
mis-stated: the algebra $\mathfrak{r}_X$ has been wrongly defined as
formed of those $u$ such that $u(X)=0$ (instead of $u(V_p)\subset X$).

\begin{corollary}[2]\label{cor:IV.A.2.4.2}
If $E$ has complex multiplication, $\mathfrak{g}_p$ is a split Cartan
subalgebra of $\End(V_p)$. If $D$ is a supplementary subspace to $X$
stable under $\Gal(\bar K/K)$, then $X$ and $D$ are the characteristic
subspaces of $\mathfrak{g}_p$ and the inertia algebra $\mathfrak{i}_p$
is the subalgebra of $\End(V_p)$ formed by those $u\in\End(V_p)$ such
that $u(D)=0$, $u(X)\subset X$.
\end{corollary}

The proof is analogous to the one of Cor.~\ref{cor:IV.A.2.4.1} (and in
fact simpler).
